\section*{Hoofdstuk \ref{ch:b3:trace-element-modelling} --- Sporenelementmodellering}

\begin{solution}{exo:b3:trace-element-modelling:1}
Spinelliet-lherzoliet: $0.12\times0.40 + 0.88\times0.016 = 0.048 + 0.0141 =
0.062$. Granaatlherzoliet: $0.195\times0.40 + 0.167\times0.11 +
0.639\times0.016 = 0.078 + 0.0184 + 0.0102 = 0.107$.
\end{solution}

\begin{solution}{exo:b3:trace-element-modelling:2}
$1/(0.05 + 0.02\times0.95) = 1/0.069 = \qty{14.5}{ppm}$;
$1/(0.05 + 0.10\times0.95) = 1/0.145 = \qty{6.9}{ppm}$.
\end{solution}

\begin{solution}{exo:b3:trace-element-modelling:3}
De verrijking is het grootst naarmate $F \to 0$: $1/D = 100$. Met $D = 0$,
$C_L/C_0 = 1/F = 20$.
\end{solution}

\begin{solution}{exo:b3:trace-element-modelling:4}
$\mathrm{La/Yb} = 30/2.5 = 12.0$; $(\mathrm{La}/\mathrm{Yb})_N =
(30/0.232)/(2.5/0.163) = 129.3/15.34 = 8.4$.
\end{solution}

\begin{solution}{exo:b3:trace-element-modelling:5}
Batch: $C_S/C_0 = 0.01/(0.01 + 0.05\times0.99) = 0.168$. Fractioneel:
$(0.95)^{99} = 0.0062$, 27 keer lager. De momentane smelt bevat
$0.0062/0.01 = 0.62\,C_0$: minder dan de bron.
\end{solution}

\begin{solution}{exo:b3:trace-element-modelling:6}
$D_0 = 0.58\times0.0003 + 0.27\times0.002 + 0.12\times0.053 + 0.03\times0.01
= 0.00737$; $P = 0.20\times0.0003 + 0.25\times0.002 + 0.45\times0.053 +
0.10\times0.01 = 0.0254$. Bij $F = 0.02$: niet-modaal $1/(0.00737 +
0.02\times0.9746) = 37.2$, modaal $1/(0.00737 + 0.02\times0.9926) = 36.7$. Bij
$F = 0.10$: 9.54 en 9.38, \qty{1.7}{\percent} uit elkaar.
\end{solution}

\begin{solution}{exo:b3:trace-element-modelling:7}
$D = 0.1$: batch $1/(0.1 + 0.09) = 5.26$; geaggregeerd $(1 - 0.9^{10})/0.1 =
6.51$. $D = 0.01$: batch $1/(0.01 + 0.099) = 9.17$; geaggregeerd
$(1 - 0.9^{100})/0.1 = 10.00$. Het paar met $D = 0.1$ verschilt meer: voor
$D = 0.01$ bevatten beide smelten al bijna het hele element, en beide naderen
$1/F$.
\end{solution}

\begin{solution}{exo:b3:trace-element-modelling:8}
$D' = 0.99\times0.001 + 0.01 = 0.01099$. Eerste smelt: $1/D' = 91$, tegenover
1000 bij \omterm{thm:b3:trace-element-modelling:fractional}{fractioneel smelten}. Bij $F = 0.01$: $91\times0.99^{90.0} = 36.8$,
tegenover $1000\times0.99^{999} = 0.044$: de poriesmelt voorkomt dat het residu
wordt leeggetrokken.
\end{solution}

\begin{solution}{exo:b3:trace-element-modelling:9}
De verhouding wordt verhoogd met $(D_{\mathrm{Sm}} + F(1 - D_{\mathrm{Sm}}))/(D_{\mathrm{La}}
+ F(1 - D_{\mathrm{La}}))$: bij $F = 0.01$, $0.0714/0.0177 = 4.0$; bij $F = 0.10$,
$0.156/0.107 = 1.46$. De verhouding is gevoelig voor $F$ bij kleine $F$ en vrijwel
ongevoelig daarboven, bij ongeveer \qty{10}{\percent}; ze hangt ook af van $D_{\mathrm{Sm}}$,
dus van de residumineralen.
\end{solution}

\begin{solution}{exo:b3:trace-element-modelling:10}
$\frac1F\int_0^F \frac1D(1 - x)^{1/D - 1}\,\mathrm dx = \frac1F\bigl[-(1 -
x)^{1/D}\bigr]_0^F = [1 - (1 - F)^{1/D}]/F$, de formule uit het hoofdstuk; als
$(1 - F)^{1/D} \to 0$, nadert ze $1/F$ maal $C_0$. Binnen \qty{1}{\percent}:
$(1 - F)^{100} < 0.01$, $F > 1 - 0.01^{0.01} = 0.045$, ongeveer
\qty{4.5}{\percent}.
\end{solution}

\begin{solution}{exo:b3:trace-element-modelling:11}
Helling $a = (10 - 6)/(20 - 10) = \qty{0.40}{ppm^{-1}}$; snijpunt $b = 10 -
0.40\times20 = 2.0$. $D_{\mathrm{Yb}} = 0.40\times0.44 = 0.176$;
$C_0^{\mathrm{La}} = 2.0\times0.44/0.824 = \qty{1.07}{ppm}$. $F_P = 1.07/20 =
\qty{5.3}{\percent}$, $F_Q = 1.07/10 = \qty{10.7}{\percent}$.
\end{solution}

\begin{solution}{exo:b3:trace-element-modelling:12}
$F = x_i/p_i$: clinopyroxeen $0.12/0.45 = 0.27$, spinel $0.03/0.10 = 0.30$,
orthopyroxeen $0.27/0.25 = 1.08$, olivijn $0.58/0.20 = 2.9$; clinopyroxeen
smelt het eerst. Bij $F = 0.2$: $\bar D = (0.00737 - 0.0254\times0.2)/0.8 =
0.0029$. Voorbij $F = 0.27$ veranderen de $p_i$ (er is geen clinopyroxeen meer
om te smelten), en de formule, die constante $p_i$ veronderstelt, geldt niet meer.
\end{solution}

\begin{solution}{pb:b3:trace-element-modelling:1}
\textbf{1.} $D_{\mathrm{La}} = 0.12\times0.058 + 0.88\times0.001 = 0.00784$;
$D_{\mathrm{Yb}} = 0.12\times0.58 + 0.88\times0.16 = 0.210$.

\textbf{2.} $D_{\mathrm{La}} = 0.195\times0.058 + 0.167\times0.001 +
0.639\times0.001 = 0.0121$; $D_{\mathrm{Yb}} = 0.195\times0.58 +
0.167\times4.5 + 0.639\times0.16 = 0.967$. Afgerond: 0.0078, 0.21, 0.012, 0.97.

\textbf{3.} $D_{\mathrm{Yb}}$, van 0.21 naar 0.97, omdat granaat de zware
zeldzame aarden vasthoudt ($D = 4.5$). Ytterbium is nog heel licht incompatibel,
$0.97 < 1$.

\textbf{4.} $(0.648/0.232)/(0.441/0.163) = 2.79/2.71 = 1.03$: een vlak patroon
van ongeveer 2.8 keer chondriet.

\textbf{5.} $0.648/(0.0078 + 0.01\times0.9922) = \qty{36.6}{ppm}$; bij
\qty{5}{\percent}, \qty{11.3}{ppm}; bij \qty{10}{\percent}, \qty{6.05}{ppm}.

\textbf{6.} $F = (0.648/9.6 - 0.0078)/0.9922 = (0.0675 - 0.0078)/0.9922 =
0.0602$, \qty{6.0}{\percent}.

\textbf{7.} $\mathrm{Yb} = 0.441/(0.21 + 0.0602\times0.79) = \qty{1.71}{ppm}$
en La/Yb $= 5.6$; lava A heeft 1.7 en 5.6: dit past bij een spinelbron.

\textbf{8.} La $= 0.0078\times9.6 = \qty{0.075}{ppm}$, Yb $= 0.21\times1.71 =
\qty{0.36}{ppm}$. Overgebleven lanthaan: $0.940\times0.075/0.648 = 0.11$, ongeveer
\qty{11}{\percent}.

\textbf{9.} $F = 0.648/9.6 = \qty{6.75}{\percent}$.

\textbf{10.} $(0.648/0.0078)(1 - F)^{127.2}$: \qty{73}{ppm} bij
\qty{0.1}{\percent}, \qty{23}{ppm} bij \qty{1}{\percent}, \qty{0.03}{ppm} bij
\qty{6}{\percent}.

\textbf{11.} $0.648\times(1 - 0.940^{128.2})/0.060 = \qty{10.8}{ppm}$.

\textbf{12.} $F = \qty{6.75}{\percent}$: bij deze $F$ is $(1 - F)^{1/D} =
\num{1.3e-4}$, dus bevat de gepoolde smelt al het lanthaan van de bron, en
geldt $\bar C_L = C_0/F$, zoals voor een volledig incompatibel element.

\textbf{13.} $0.648\times0.9325^{127.2} = \qty{9e-5}{ppm}$, ongeveer 800 keer
minder dan de \qty{0.075}{ppm} van \omterm{thm:b3:trace-element-modelling:batch}{batchsmelten}.

\textbf{14.} De twee modellen geven de grenzen van het antwoord: \qty{6.0}{\percent} (batch)
tot \qty{6.75}{\percent} (gepoolde fractionele). Lanthaan legt $F$ vast binnen
ongeveer \qty{1}{\percent} smelten, wat het model ook is.

\textbf{15.} A: La/Yb $= 5.6$, $(\mathrm{La}/\mathrm{Yb})_N =
(9.6/0.232)/(1.7/0.163) = 4.0$. B: La/Yb $= 34.7$, $(\mathrm{La}/\mathrm{Yb})_N
= (15.6/0.232)/(0.45/0.163) = 24.4$.

\textbf{16.} Voor $F \to 0$: $(0.648/0.0078)/(0.441/0.21) = 83.1/2.10 = 39.6$.

\textbf{17.} $F = \qty{0.11}{\percent}$, waarbij de smelt \qty{72}{ppm} La en
\qty{2.1}{ppm} Yb zou bevatten. Lava B heeft \qty{15.6}{ppm} La: een
spinelbron kan dat niet maken.

\textbf{18.} $F = (0.648/15.6 - 0.012)/0.988 = 0.0299$, \qty{3.0}{\percent};
$\mathrm{Yb} = 0.441/(0.97 + 0.0299\times0.03) = \qty{0.454}{ppm}$, zoals
waargenomen.

\textbf{19.} Granaat vervangt spinel bij ongeveer 60 tot \qty{90}{km}: B begon
dieper dan dat te smelten, A ondieper, in het spinelveld. Een granaat-lherzolietsmelt
met \qty{9.6}{ppm} La ($F \approx \qty{5.6}{\percent}$) zou een La/Yb nabij
21 hebben, niet 5.6: A kwam dus niet uit de granaat-lherzoliet.

\textbf{20.} $a = 0.97/0.441 = \qty{2.20}{ppm^{-1}}$, $b = 0.648\times0.03/0.441
= 0.044$; bij La $= 15.6$ geldt La/Yb $= 2.20\times15.6 + 0.044 = 34.4$, tegen
34.7 voor lava B.

\textbf{21.} Olivijn verwerpt beide elementen: La/Yb blijft ongewijzigd. De
primaire smelt bevatte $9.6\times0.9 = \qty{8.64}{ppm}$ La, dus $F = (0.648/8.64 -
0.0078)/0.9922 = \qty{6.8}{\percent}$ in plaats van 6.0.

\textbf{22.} Eén bron kan beide lava's maken: B is een kleine smelt
(\qty{3}{\percent}) die begon in het granaatveld, dieper dan ongeveer 60 tot
\qty{90}{km}; A is een grotere smelt in het spinelveld.

\textbf{23.} Lava A is een \textbf{\qty[text-series-to-math=true]{6.0}{\percent}}
batchsmelt van een spinel-lherzoliet (\qty{6.75}{\percent} als de smelt
fractioneel was onttrokken en gepoold).
\end{solution}
